\documentclass[UTF8,11pt,a4paper]{ctexart}
\usepackage[margin=24mm]{geometry}
\usepackage{booktabs,tabularx,array}
\usepackage{hyperref}
\hypersetup{hidelinks,pdftitle={SRC 到 BigRIPS F2/F3 的束流强度控制}}
\setlength{\parindent}{2em}
\setlength{\parskip}{0.4em}
\newcolumntype{Y}{>{\raggedright\arraybackslash}X}

\title{SRC 出口至 BigRIPS F2/F3 的束流强度控制}
\author{DPOLAR 束线资料记录}
\date{\today}

\begin{document}
\maketitle

\begin{abstract}
本文记录从 SRC 出口到 BigRIPS F2/F3 之间可影响束流强度的装置、它们的物理作用，以及偏振氘束实验需要向 RIKEN 束线组确认的运行边界。结论先行：公开资料明确给出的 F3 总束流率上限为 $10^7$ particles/s；F1、F2 和 F2--F3（dF2）狭缝可以通过减小动量、空间或角度接受度来降低传输率，而一次束模式的安全限流主要由 SRC 上游的 beam attenuator（ATT）和联锁系统管理。
\end{abstract}

\section{束线位置与功能}

BigRIPS 第一段从生产靶 F0 到 F2，第二段从 F3 到 F7；F2 与 F3 之间是约 8.8 m 的匹配段。标准布局为
\[
\mathrm{SRC}\;\longrightarrow\;\mathrm{T\ course}\;\longrightarrow\;F0\;\longrightarrow\;D1/F1\;\longrightarrow\;F2\;\longrightarrow\;dF2\;\longrightarrow\;F3.
\]
F0 后的 D1 间隙内设置高功率 beam dump，用于截停未反应的一次重离子束；第一段的 F1 是动量色散焦点，F2 是消色差焦点，F3 是第二段的起点。

\section{可以降低强度的位置}

\begin{tabularx}{\linewidth}{@{}p{0.18\linewidth}p{0.27\linewidth}Y@{}}
\toprule
位置 & 装置 & 对束流的作用 \\
\midrule
SRC 上游低能束线 & Beam attenuator (ATT) & 插入不同衰减倍率的衰减片，降低进入 SRC--BigRIPS 的一次束强度；其插入状态被 RIBF beam interlock system（BIS）监视。\\
F0 后、D1 内 & High-power beam dump & 截停未反应的一次束，主要承担功率和辐射屏蔽任务；不是精细的连续衰减器。\\
F1 & 水平狭缝、Joker slit、degrader、beam stopper & 在动量色散平面切掉动量尾部、杂质和束晕；收窄狭缝会降低传输率。\\
F2 & 水平/垂直四维狭缝、collimator、beam stopper & 限制横向相空间和污染束流；公开 F2 配置明确列出 H/V 狭缝及 45 mm Cu beam stopper。\\
dF2（F2--F3） & 水平/垂直角度选择狭缝 & 在 point-to-parallel 光学位置把 F2 聚焦束的角度映射为空间位置，从而选择角度并降低到达 F3 的束流。\\
\bottomrule
\end{tabularx}

\section{F3 的强度限制}

RIKEN BigRIPS 官方束流说明要求在申请书中填写 ``Total Beam Rates @ F3''，并指出由于辐射法规，BigRIPS 中的总束流率不能超过 $10^7$/s。该总数包括目标核素和杂质粒子。官方建议在 LISE++ 中在 F3 后插入 Faraday Cup 2 检查该数值。

这意味着：若实验设备位于 F3 下游，必须先满足 F3 的总率限制；只在更下游用探测器触发降采样，不能替代 F3 位置的辐射和束损限制。

\section{一次束与 RI 束模式的区别}

对于 F0 产生 RI 束的标准模式，F0、D1 beam dump、F1 degrader/slit 和 F2 collimator 共同完成一次束截停及 RI 束选择。对于将一次束直接送到下游的模式，RIKEN 的 BIS 资料说明：系统根据 ATT 插入状态判断衰减是否足够；衰减不足时，会在 SRC 出口插入 Faraday cup，并可触发 beam chopper。

因此，``SRC 后保持高流强，在 F2 或 dF2 再降到允许值'' 不能仅由狭缝几何推断为可运行方案。必须确认该模式下的联锁阈值、束损位置、F2/dF2 狭缝的热负荷和屏蔽条件。

\section{对 after-SRC 偏振仪的建议}

偏振氘束实验建议把需求拆成两个独立量：
\begin{enumerate}
  \item after-SRC 极化计处的一次束强度和允许的探测器计数率；
  \item F3 处（若束流继续进入 BigRIPS 下游）的总束流率，包括杂质。
\end{enumerate}
优先向束线组询问是否存在专门的局部衰减器或可接受的 F2/dF2 准直方案，而不要把普通 F2 狭缝直接当作高功率 beam attenuator。若需要保持偏振，任何材料衰减器还应单独评估多重散射、能损、核反应和可能的自旋系统误差。

\section{资料来源与目录组织}

本文件放在\texttt{docs/beamline\_info/bigrips/}，与已有资料分工如下：
\begin{itemize}
  \item \texttt{overview.md}：BigRIPS 光学分段、焦点和基本参数；
  \item \texttt{primary-beam-and-polarized-deuteron.md}：一次束和偏振氘束的能量、流强及加速链；
  \item \texttt{slits/acceptance-and-loss.md}：狭缝接受度、束流损失和纯度权衡；
  \item 本文件：SRC--F3 的限流位置、F3 率限制和一次束联锁问题。
\end{itemize}

\begin{thebibliography}{9}
\bibitem{bigrips-config} RIKEN Nishina Center, ``BigRIPS Configuration'',\newline
\url{https://www.nishina.riken.jp/ribf/BigRIPS/config.html}.
\bibitem{bigrips-intensity} RIKEN Nishina Center, ``Secondary beam intensity expected'',\newline
\url{https://www.nishina.riken.jp/ribf/BigRIPS/intensity.html}.
\bibitem{bigrips-focal} BigRIPS Team, ``F1/F2 chamber technical information'',\newline
\url{https://ribf.riken.jp/BigRIPSInfo/chamber/f1.html},\newline
\url{https://ribf.riken.jp/BigRIPSInfo/chamber/f2.html}.
\bibitem{f2-f3-slit} K. Yoshida et al., ``Slit system between the foci F2 and F3 of the BigRIPS separator'', RIKEN Accel. Prog. Rep. 52 (2019), p. 130,\newline
\url{https://www.nishina.riken.jp/researcher/APR/APR052/pdf/130.pdf}.
\bibitem{bis} M. Komiyama et al., ``Development of a machine protection beam interlock system for the RIBF'', RIKEN Accel. Prog. Rep. 58 (2025), pp. 98--99,\newline
\url{https://www.nishina.riken.jp/researcher/APR/APR058/pdf/98.pdf}.
\end{thebibliography}

\end{document}
